\section*{Proposed Solution} \label{sec:solution}
\noindent

We introduce [ACRONYM], a novel approximate error correction scheme that replaces algebraic coding with a hash-constrained two-dimensional grid structure and a constraint-guided combinatorial search decoder. [ACRONYM] is designed around three core ideas: \emph{(i)} scatter errors uniformly before hashing, \emph{(ii)} create localized, overlapping constraints that jointly identify error locations, and \emph{(iii)} prune the search space using already-correct constraints so that exhaustive enumeration remains tractable even at high error rates.

\textbf{Encoding.} Given a data block of $L$ bits, [ACRONYM] first applies a keyed, invertible Feistel permutation to the bit indices, reordering the bits pseudorandomly. The permutation uses a multi-round Feistel network with SHA-256 as the round function and employs the cycle-walking technique to handle domain sizes that are not powers of two, guaranteeing a bijection on $\{0, \ldots, L-1\}$ for any $L$. The permuted bits are then arranged into a square grid of $N \times N$ cells, where $N = \lceil \sqrt{L} \rceil$, padded with zeros if necessary. CRC digest values are computed over each row and each column of the grid and stored alongside the data as error-correction metadata. The total metadata is $2Nh$ bits, where $h$ is the CRC width; for $L = 4{,}096$ bits and $h = 32$ (CRC-32), this is exactly 4,096 bits---100\% overhead---and scales as $O(1/\sqrt{L})$.

\textbf{Decoding.} Upon reading a block that may contain errors, [ACRONYM] reconstructs the same grid via the same Feistel permutation and recomputes row and column CRC values. Rows and columns whose recomputed digests differ from stored values are marked as \emph{mismatched}; the rest are \emph{matched}. A constraint graph is then built: nodes are hash digest computations (rows and columns) and edges connect any two nodes that share bit coverage---in the default single-row/column configuration, this is a complete bipartite graph where each edge points to exactly one grid cell. The decoder then executes a greedy, constraint-guided bit-flip search. Mismatched nodes are processed in ascending order of covered-bit count. For each mismatched node, the search space is pruned by the \emph{matched-neighbor pinning} heuristic: bits already covered by correctly-verified neighboring nodes are excluded (``pinned''), reducing the combinatorial search from $\binom{N}{k}$ to $\binom{|\text{candidates}|}{k}$ at each flip level $k$. The decoder commits the flip combination that maximizes total matched nodes globally, updates node labels, and repeats until all nodes are matched or no further progress is achievable.

\textbf{Key result.} At 4,096 bits and 5\% BER (200 flips), [ACRONYM] with CRC-32 achieves $\geq$95\% full correction at 100\% overhead, versus BCH's $\sim$174\% overhead for equivalent capability. The Feistel permutation converts burst errors into the equivalent of random errors, so [ACRONYM] achieves statistically identical correction rates for random and burst error patterns without any additional mechanism. The decoder operates entirely with standard CRC computations---no finite-field arithmetic is required at any stage.
